\section{Round-ten reconstruction: a past Markov kernel, an eigenfunction Doob transform, and exact memory}
\label{sec:r10-a4}

Let \(\mathsf H\) be the weighted natural-extension past space.  A point of
\(\mathsf H\) contains only coordinates at times \(s\le0\); it does not
contain the unstable/future itinerary.  Disintegration of the prepared
invariant path law over this past sigma-field gives a regular conditional
future kernel \(P_t(h,dh')\).

\subsection{The past Feller process}

\begin{theorem}[Weighted past-kernel Feller semigroup]
\label{thm:r10-a4-feller}
The conditional kernels \((P_t)_{t\ge0}\) form a conservative strongly
continuous Feller semigroup on a weighted uniformly continuous history space.
For histories agreeing on the last \(L\) units,
\[
 |P_tF(h)-P_tF(\tilde h)|
 \le C\|F\|_{\rm dyn}(e^{-\gamma L}+e^{-\gamma t})
\]
after the singularity-shield error is included.  The future law is
nondegenerate unless the prepared phase itself is deterministic.
\end{theorem}

\begin{proof}
The natural extension disintegrates a Gibbs path law into a one-sided
\(g\)-kernel in the forward temporal orientation.  Summable variations give
uniqueness and exponential coupling of two futures whose finite pasts agree.
The A3 singularity shield transfers the symbolic coupling to graph-completed
billiard paths, while the exponential roof tail controls continuous-time
renewal residuals.  Conditioning successively and concatenating futures gives
Chapman--Kolmogorov.  The coupling estimate gives Feller continuity and strong
continuity.  Since the future coordinate was not retained in \(h\), the
conditional kernel is not a Dirac kernel.
\end{proof}

\subsection{The correctly oriented Poisson decomposition}

Let \(g\) be centered and dynamically H\"older, and write
\(Pg(h)=\int g(h')P_1(h,dh')\).  Put
\[
 h_g=\sum_{n\ge1}P^ng,
 \qquad
 m_g(h_0,h_1)=g(h_1)+h_g(h_1)-Ph_g(h_0).
\]

\begin{lemma}[Conditional martingale--coboundary decomposition]
\label{lem:r10-a4-poisson}
The series for \(h_g\) converges in the weighted H\"older norm and
\[
 \mathbb E[m_g(H_k,H_{k+1})\mid\mathcal F_k]=0.
\]
Moreover
\[
 \sum_{j=1}^n g(H_j)
 =\sum_{j=0}^{n-1}m_g(H_j,H_{j+1})
   +Ph_g(H_0)-Ph_g(H_n).
\]
\end{lemma}

\begin{proof}
The Feller coupling estimate implies
\(\|P^ng\|\le Ce^{-\gamma n}\|g\|\), so the Poisson series converges.  By the
Markov property,
\[
 \mathbb E[g(H_{k+1})+h_g(H_{k+1})\mid\mathcal F_k]
 =P(g+h_g)(H_k)=Ph_g(H_k),
\]
because \((I-P)h_g=Pg\).  This proves conditional centering.  Summing the
definition and using the Poisson identity gives the coboundary formula.  The
raw centered observable is never called a martingale difference.
\end{proof}

\begin{theorem}[Uniform quenched enhanced invariance principle]
\label{thm:r10-a4-quenched}
For every compact regular family of prepared phases and every vector of
weighted dynamically H\"older observables, the conditional laws given
\(H_0=h\) of the diffusively rescaled sums and their canonical iterated
integrals converge, locally uniformly in \(h\), to a Brownian rough path.  Its
covariance and area drift are the predictable bracket and antisymmetric
coboundary limits of Lemma~\ref{lem:r10-a4-poisson}.
\end{theorem}

\begin{proof}
The martingale increments have a uniform \(2+\delta\) moment under the
conditional kernels, by the roof and singularity exponential moments.
Conditional Lindeberg and bracket convergence follow from Feller coupling and
the ergodic theorem uniformly on compact history sets.  The martingale rough
functional CLT gives the first two levels.  The coboundary is negligible at
the first level and contributes the explicit antisymmetric area drift at the
second.  Truncation at Young returns and the A3 exponential shield remove the
compact-history restriction.  This proves stable, hence quenched, convergence.
\end{proof}

\subsection{The correct nonlinear history semigroup}

For a bounded potential \(\Psi\), let
\[
 T_t^\Psi F(h)=\mathbb E_h\left[
 e^{\int_0^t\Psi(H_s)ds}F(H_t)\right].
\]
Let \(e^{\lambda_\Psi t}\), \(h_\Psi>0\), and \(\ell_\Psi\) be the simple
leading Feynman--Kac eigen-data on the regular source chart.

\begin{theorem}[Eigenfunction-normalized nonlinear semigroup]
\label{thm:r10-a4-doob}
The Doob kernels
\[
 \widetilde P_t^\Psi F
 =e^{-\lambda_\Psi t}h_\Psi^{-1}T_t^\Psi(h_\Psi F)
\]
are conservative and satisfy the exact semigroup law.  Consequently
\[
 \mathcal U_t^\Psi F=\log\widetilde P_t^\Psi(e^F)
\]
is a cash-additive nonlinear semigroup.  The unnormalized logarithm
\(\log T_t^\Psi(e^F)\) also composes; the state-dependent quotient
\(\log T_t^\Psi e^F-\log T_t^\Psi1\) is not used as a semigroup unless it is
expressed through the displayed Doob transform.
\end{theorem}

\begin{proof}
The Feynman--Kac operators compose by the Markov property and satisfy
\(T_t^\Psi h_\Psi=e^{\lambda_\Psi t}h_\Psi\).  Substitution therefore gives
\(\widetilde P_{t+s}^\Psi=
\widetilde P_t^\Psi\widetilde P_s^\Psi\) and
\(\widetilde P_t^\Psi1=1\).  Exponentiating and taking logarithms gives the
nonlinear tower and cash additivity.
\end{proof}

\subsection{Exact compressed-resolvent memory}

Let \(L_\Psi\) be the generator of the Doob phase on its common Hilbert
realization, \(P\) a finite-rank orthogonal resolved projection, and
\(Q=I-P\).  Put
\[
 C(z)=P(z-L_\Psi)^{-1}P,
 \qquad
 \widehat K(z)=zP-PL_\Psi P-C(z)^{-1}.
\]

\begin{theorem}[Full-Q dilation and causal memory]
\label{thm:r10-a4-memory}
On \(P\mathcal H\oplus Q\mathcal H\), the fixed operator
\[
 \mathbb L_\Psi=
 \begin{pmatrix}
 PL_\Psi P&PL_\Psi Q\\
 QL_\Psi P&QL_\Psi Q
 \end{pmatrix}
\]
has compressed resolvent \(C(z)\).  The resolved correlation satisfies
\[
 C'(t)=PL_\Psi P\,C(t)+(K*C)(t),
 \qquad C(0)=P,
\]
with the displayed sign convention.  After adjoining the finite principal
parts of \(C(z)^{-1}\) as descriptor modes while retaining the full
\(Q\)-space, the residual kernel is exponentially integrable on the spectral
phase strip.
\end{theorem}

\begin{proof}
The Schur complement of \(z-\mathbb L_\Psi\) is
\[
 zP-PL_\Psi P-PL_\Psi Q(z-QL_\Psi Q)^{-1}QL_\Psi P=C(z)^{-1}.
\]
Rearranging gives the stated \(\widehat K\).  Inverse Laplace transformation
yields the Volterra equation and its initial value.  Meromorphic continuation
from the A2 spectral packet has finitely many principal parts in every closed
strip.  A finite descriptor realization represents exactly those parts; it
is coupled to, not substituted for, the unresolved \(Q\)-space.  After their
removal the contour moves to a negative line, giving exponential integrability
of the residual memory.
\end{proof}

\begin{theorem}[History pressure and rough tangent]
\label{thm:r10-a4-main}
The eigenfunction-normalized history semigroups converge under the diffusive
scaling of Theorem~\ref{thm:r10-a4-quenched} to the exponential semigroup of
the Brownian rough differential equation with the proved covariance and area
anomaly.  The resolved memory coordinates converge jointly to the Volterra
Gaussian lift generated by Theorem~\ref{thm:r10-a4-memory}.
\end{theorem}

\begin{proof}
Quenched rough convergence and uniform exponential moments pass bounded
cylinder Feynman--Kac functionals to the rough limit.  Analytic perturbation
passes \(\lambda_\Psi\) and \(h_\Psi\) through the scaling, so the Doob
normalization converges as well.  The exponentially integrable residual
memory is a continuous convolution functional, and the finite descriptor
modes are finite-dimensional linear functionals.  Joint convergence follows
from the continuous mapping theorem.
\end{proof}
